\chapter{Delight}

\vspace{-1.3cm}
\begin{localsize}{10}
	\begin{quote}
		``If a man does not keep pace with his companions, perhaps it is because he hears a different drummer. Let him step to the music he hears, however measured or far away."
		\begin{flushright}- Henry David Thoreau \end{flushright}
	\end{quote} 
\end{localsize}
\vspace{1cm}

The sun slowly rose in the eastern sky; innocent rays of morning trickled through the dense forest canopy.
The sparse beams reached for the undergrowth, fighting off the morning mist like a potent antidote.
The daily struggle between dew and vapor began anew.
Kraerion sat awake in his hollow, hunched on his forepaws, gazing contentedly out through the morning air.
He had awaken---like most mornings---to the varied orchestra of nightingales and blackbirds, tirelessly competing for the vibrations of the humid air.
The beautiful song was now and then disrupted by the distinct cooing of one of the neighbouring wood pigeons.

Kraerion had managed well: gathering his own feed, supplemented by the odd job here and there when such opportunities arose.
But---of course---he was less efficient now than when he was still employed.
Yet he enjoyed his new, savage life; what he lost in feed and comfort, he gained in beauty.
The absence of the natural irregularities of life---which a steady job, with an even steadier wage, had smoothed numb---was to him its own kind of poverty.
And waking before the sun instead of with it---there was a cost to be paid there, too.

\renewcommand*{\thepage}{\footnotesize \hexadecimal{page}}

Before he left his burrow and former life, he had saved for a long time.
He had managed to obtain his small fortune of seven screw-nuts by sheer hard work: three for the birch, one for Redrill, and the three remaining to act as a safety net during his transition to independence, to cover any inherent uncertainty that his new life would surely entail.

But that summer, he was blessed with something magical: prices began to fall.
Little by little, and then all at once.
By the time the sun reached its summer solstice, the prices had dropped to about half of what they had been during the spring---a drop that made his savings ample not only for his transition, but for anything the future could possibly bring.

No cause could be attributed to the sudden drop in prices, but neither thought nor concern burdened the Wadborough animals; instead, they feasted as if the abundance of summer would never be depleted.% break

The exchange to leaves had happened exactly as Eaglewing had dictated, one week after the meeting in the glade.
The animals had formed a queue before the Old Oak to attempt an orderly, civilized proceeding.
Although the animals had increased in affluence, queues were still as innate to them as a sheep stalking a pack of wolves, or a chicken snagging a fox for dinner; thus, it was nothing short of a miracle that no major incident occurred.
Kraerion, ever cautious, had waited until the absolute end to avoid most of the expected tumult before exchanging his remaining screw-nuts for a thin stack of leaves.\\

The sound of scratching against the bark broke his peace. % Peace, contemplation, thoughts
As so often happened to Kraerion, his thoughts remained but half-processed---he was flung back to reality leaving him momentarily disoriented.
He leaned his head out of his hollow to peer down the trunk and saw a familiar shape standing at the foot of his birch.

\renewcommand*{\thepage}{\footnotesize \arabic{page}}

``I'm not climbing this \textit{thing}!" Cinnamon cried mockingly.
In response, Kraerion moved down the tree and stopped roughly a meter from the ground, in a manner not dissimilar to that of a nuthatch.

``You never said goodbye," she said at length when Kraerion answered only with silence.

``What could I have said?" he answered, knowing full well how stupid it sounded. ``Besides, I don't think I could---it was difficult enough as it was."

``Well, if it was too difficult, might it not be the wrong decision?"

``No!" He cried harshly as if by reflex, ``It was truly the best decision I've ever made."

As usual, he spoke before thinking and failed to consider the impact of his words---now Cinnamon looked rather hurt.

``I miss you," he added quickly.
``You know what I mean.
You know how miserable I was; I wouldn't have been good for you, or anyone else, had I stayed."

Cinnamon stared down at her paws, kicking at a small pinecone lying next to her as she processed his words.

``I guess you're right..." she whispered.

After a moment of uncomfortable silence, Cinnamon finally met his gaze; the hurt hadn't entirely vanished from her eyes.

``Do you want to walk somewhere?" he asked, more to get them out of the awkward silence than anything else.

``Somewhere?"

``Anywhere."

``Didn't you always lecture me about preferring to wander?" She giggled.
``That walking was for drones without a soul?" She found his eyes again, making something that could be construed as a smile before turning.
Before Kraerion had a chance to reply, she started walking east, and he was left without a choice save to catch up with her.

Kraerion had long made an effort to distinguish between walking and wandering.
``Walking is active and with purpose---you \textit{walk} with the purpose to get somewhere," he would preach.
``But wandering \textit{is} the purpose, done for its own sake---an end in itself."
To Kraerion, they were more than just definitions; they symbolized the two different outlooks of life: the former, the ambitious numbness of habit; the latter, the very beauty of life itself.
He reasoned that one's quality of life could be deduced from this innocent choice of words.
Someone who walked instead of wandered, he believed, would eventually meet with the Thoreauvian Fear: to wake up one day and discover that they had not lived.
The chipmunks---and thus also Cinnamon---had built their existence around transportation, schedules, and habit, so whilst she would always listen, he feared she would never truly hear.\\

After wandering for some time, Cinnamon broke the silence.

``I'm worried about your sister," she began tentatively.
``She's changed."

``Veritas?" Kraerion asked.

``I--- I don't know. Well, maybe not changed per se.
She's herself alright, only more so."

``Honesty, can you have too much of that?" Kraerion chuckled.

``Kraerion!
You're not listening."

Kraerion paused, giving it some more thought, "I think I see what you mean, though.
She can get stuck in her head sometimes, not letting thoughts go---like being caught in a never-ending loop.
But she's sharp; I'm sure she'll figure it out."

``I don't know, maybe..." She said slowly. ``Maybe that's why we call someone smart \textit{sharp}---they cut wherever they turn; themselves too if not supervised.
Whilst someone dull couldn't hurt a fly if they tried.
But you know what, Kraerion? you were always her whetstone, always there to keep her on track.
And well, you're not there anymore..."

``Nonsense, you talk as if she's crazy!"

``Maybe she is crazy! She's at least beginning to be perceived as if she is," Cinnamon said, not backing down.
``But, there's more: she's started talking about our economy as if there's some great danger to it.
She's not terribly coherent in the best of times, but now she's regressed to repeating---repeating these incomprehensible phrases."

``What do you mean phrases---what phrases?"

``There are a bunch of them, but there are a couple in particular that she keeps repeating which I managed to remember: 'Feed is dialectic; hunger is constant,' and 'The delta must eventually reverse,'" she said. ``Do they tell you anything?"

``Well..." He thought for a while. ``The first is clearly about winter. Feed is cyclical: from the first rare morels that nigh beat the thaw, to the late autumn tubers that are accessible until the ground freezes.
And hunger is... well, hunger.
It's constant and independent of season.
But why would she bring it up now?
That's always been the case, and the Winter Fund was invented to fix it.
She's \textit{practically} running the Fund now, in all but name." He took another pause.
``The second one I don't understand, though."

Their conversation did not end, but transposed to one in silence.
With words omitted, they traversed through the forest.
Animals communicate more with their bodies than humans, and the two chipmunks soon fell into the comfort they had felt so many times before in each other's presence.
Nevertheless, there was much left unsaid between them, and a tension slowly grew.
Kraerion suspected she had not shown up just for her sense of duty: to tell of Veritas's well-being.
He was certain she felt it too, yet they left it unresolved.


She stopped abruptly, as if she could not bear the silence a second longer.
``You are never coming back, are you, Kraerion?" she asked, her two chipmunk eyes---glimmering like black pearls caught in sunlight---glaring straight through him.

``No," he said staunchly after some thought.
``I'm sorry, but I don't think I ever will."

``But why---why would you leave?
I'm not asking for Veritas, or Aequitius, or anyone else; I'm asking for me, Kraerion.
For me.
What about us---what about the life we were creating?"

``I don't know what more I can say," he said sadly. ``I can explain how awful I find life in darkness with that damnable filthy soil all around, how it's like my body completely rejects it, just as you seem to reject the notion of living in the trees.
No matter how much I want to be with you, I can never be in consonance with a life underground.
I assure you, it hurts me as much as I think it hurts you to not see you every day."

Once again, they regressed to silence.
Exhausted by their imposition, they wandered along.

They wandered past the rabbits' great hill, and thence through the clearings which constitute the central parts of the forest.
Many an animal lay leisurely basking in the morning sun, watching as the two chipmunks quietly passed by.
``That's queer," Kraerion thought.
``Aren't they supposed to be working at this hour?"
But he soon dismissed the thought.\\

When they reached the Glade of Representatives, Kraerion could scarcely recognize it without the bustling presence of the forest's animals.
The glade, however, wasn't empty.
On the contrary, it was in a frenzy.
Woodpeckers swarmed the Old Oak---pecking away at its trunk.
Sturdy roebucks hauled loads of dirt away from the center, and crows were all around, talking amongst themselves---making plans.

They had nigh entered the glade before a voice called out for them.

``Good morning!" cried a voice, which they now saw belonged to Hroðulf, as he waved them over.
He stood next to a straight stick planted perpendicular to the ground; stones had been laid out around it to form a semicircle.
In between every stone lay a couple of pebbles.

``Morning," Cinnamon responded as they approached.
``What is this?" she continued, pointing at the stick and stones with her front paw.
Hroðulf only smiled for a while, then he spoke:

\input{second_version/poem}

Just when he stopped talking---as by a bespoke request to God---the sun appeared from behind the lone cloud in the sky.
Its rays gave birth to a shadow running from the stick to the second westernmost stone.

``Wait---" Cinnamon said, astonished.
``Do you mean that this thing---this \textit{contraption}---tells us the time?"

``Yes!" the crow cawed with delight.
``And here I thought \textit{rodents} were completely void of wit!
Whilst the sun travels above us from east to west, its shadow travels from west to east.
The stones simply enumerate the time of day."

``And how accurate is this shadow?" Cinnamon asked, ignoring the slight.
``Does it fall on the same stone at the same time every day?"

``Not only is it accurate, it's also completely independent of location.
No matter where in the Forest you are, the angle of the shadow is exactly the same.
We can just raise the same structure in all parts of our Forest, which I presume you chipmunks---with your vocation---will see the value in."

``This is truly \textit{brilliant}," Cinnamon exclaimed, her eyes following the shadow.
``Nothing will be the same again!"

``No, this is nothing.
It's simply a natural, iterative step of progress," Hroðulf concluded, albeit with a tone free from humility.
His smile from before appeared again.
Kraerion watched him hesitate, as if he considered the prudence of sharing his inner thoughts.
Clearly, his urge won in the end, because he continued:\\

``But I will tell you what is: sentiment.
There are so many of these silly beliefs that, by themselves, become the agents of their own truth.
For us as individuals, spending less than we earn is the only viable strategy for long-term wealth; hence, it is assumed that the same strategy ought to be applied at the aggregate---to the total spending of the Forest.
But naught could be further from the truth!
What is true for the individual might be completely false for the Forest as a whole.

``Consider, for example: if each Animal saves more, then each business would receive less.
Subsequently, they would have less to pay out in wages to the Animals they employ, which in turn would only decrease what each Animal has to spend in the first place.
Good individual decisions to save could cause a vicious spiral of less and less spending, weakening our Forest.
Frugality would only make us all poorer!

``The savings of every Animal bear no meaning when you sum them together, if you really think about it.
If an Animal reduces her savings by way of consumption, those leaves only become the savings of someone else.
The total savings of the Forest must always equal the total amount of leaves.

``And that leaves us with sentiment.
It has been shown that Animals save less if they believe in a prosperous future, and save more if they believe in a bleak one.
How backwards is it that our Forest is at the whim of fools' belief!
And what if I told you there might be other ways for us to control the total spending of the Forest?"\\

When Cinnamon responded to Hroðulf's conjectures, Kraerion had long since stopped listening.
He didn't care for words like spending and savings, nor for the forest at large.
But the sundial bothered him greatly.
``The need to know the time \textit{is} the problem," he muttered with disgust.
``No one makes plans to smell the hillside flowers, nor watch the glimmer of the morning dew; when the day's activities become shackled to time, the wonders of the world do not cease to exist, but become imperceptible when there is no time left to stumble upon them.
The solution to Time is indeed no solution at all."

Instead, the chipmunk admired the woodpeckers' hard work.
He was hit by a sudden realization: weren't there too many woodpeckers for just one tree?
And the roebucks were hauling dirt, but there was no sign of digging---whence came the dirt?
But then again, what did Kraerion know about the process of industry---a subject he had so eagerly chosen to ignore?\\

\renewcommand*{\thepage}{\footnotesize \mayadigit{42}}

When great ideas are discovered---those seemingly obvious in hindsight---they appear imperishable.
There's no way to put the genie back in the bottle.
They propagate from mind to mind---like a virus---only by the perceived virtue of the ideas themselves, until they're safely kept in the collective mind of society.
Yet great ideas do perish; they slowly change, transform, and morph until they are something completely new, and nothing of their original beauty remains.

We humans of today have learned to tell time in base 60 because long ago, when the ancient Babylonians created their first sundial, they had in turn inherited a sexagesimal system from the even older Sumer.
Much of that meaning---the innate substance of their mathematics---has long been forgotten,or fallen out of relevance, or transposed into what we consider maths now.
Yet, to this day, the arbitrary number 60, written in decimal, remains inert.
Now is but the cumulative residue of legacy; the fabric of time woven with fragments of lost meaning.

Or is anything, really, arbitrary?
Since 60 is a fortunate number for mathematics, maybe the Babylonians discovered how to measure time \textit{because} of their inherited numbers.
Their discovery was but the product of the trajectory of their past---propelling them to take the next progressive step of human understanding.
Could Newton or Leibniz have developed calculus if they were burdened with doing their long division with Roman numerals?

\renewcommand*{\thepage}{\footnotesize \arabic{page}}

The survival of ideas has by some been associated with Darwin's Natural Selection: knowledge---that which is useful to its carrier (or not so useful, predatory even, as with the spreading of a cult)---is remembered and passed on.
Information that is easily shared, like an infectious joke or a compelling fireside tale, survives the erosion of time.
And just like in biology, numbers like 60, like our useless appendix, still remain inert---a relic of the past without a function except for the telling of time.
The phenomenon might best be exemplified with languages; they slowly change, become dialects, merge and borrow words, some dying whilst others spread to conquer the earth!
But as with any stochastic process, these changes appear---insofar as we are capable of discerning them---as noise; and only when we zoom out, over the vast time-frames that are our history, are we able to see the undeniable signal.

Whatever our history has to teach, when Hroðulf saw a shadow and realised what it meant, he shackled the future to his preference.
Did he consider the future crows, who would grow up burdened with learning his clock---that the arbitrary number of stones he chose for his sundial would be the last piece to remain of his existence on earth?
Did he consider the implications of his societal economic experiments, what society they would create, and that they, once set loose, would be unstoppable and irreversible when reaching motion?\\

By the time Cinnamon finally exhausted her curiosity, the sun was about to go down.
The two chipmunks began their journey back to Kraerion's hollow.
They took a different route back, to parts of the forest other than those they had seen on their way there.
They wandered once again in silence, but this time it was different and not burdened with tension.
It was their silence, the one they had shared for so many years; a silence to grow old in.
Both felt it, yet they knew that their paths would soon diverge.

``You know, it's not that far between your burrow and my hollow..." Kraerion began cautiously.

``I know," Cinnamon said slowly.

``You think we could give it a try?"

``I don't know, maybe..." she almost whispered, but then she nodded.
``Sure."
They didn't exchange another word the short distance back to the hollow; they didn't need to.
Kraerion thought back to Hroðulf's line, 'our sun, eternal bound', thinking of the sun in his life and entertaining the idea; he smiled a solemn smile at his own folly.
